% Human and machine writing, and the Discussion. Numbers come from numbers.tex (generated).

\subsection{Human and machine writing}

A measure that is cheap and does not depend on length can be pointed at a current question: does text written by a language model carry feeling differently from text that people write on the same prompt? Earlier comparisons mostly report that model output is less negative, and more neutral or positive, than human text\cite{guo2023,kabir2024,munozortiz2024,tian2024}, with exceptions by domain\cite{ha2024,sui2026}, alongside differences in vocabulary and grammar\cite{reinhart2025,kobak2025}. Before generating or scoring any machine text I fixed ten length-free measures read from signatures and four predictions (Methods): machine text would be more positive overall (HL1), use fewer negative words (HL2), spread its charge more evenly over sentences (HL3), and close on a more positive note than its body (HL4).

\figMachine

Two comparisons were made. HC3 pairs questions with an answer from a person and an answer from ChatGPT\cite{guo2023}; \hcN{} pairs had answers long enough to score. For essays, I sampled \esN{} essays that US school students wrote on eight assignments\cite{crossley2024} and had four Claude models write to each student's assignment and length, once on the bare assignment and once with the instruction to write as a student of that grade would. The writing agents were told nothing about the study.

In HC3 all four predictions hold, and all four effects are small: ChatGPT answers are higher in net density by \hcNet{} human standard deviations, lower in dolor density by \hcDolors, lower in clustering by \hcClustering{} and higher in ending lift by \hcEndingLift{} (\figref{fig:machine}a). In the essays only the last holds under both instructions. The final fifth of a machine essay sits \esLiftMachineA{} per thousand words above the essay's own density, against \esLiftHumanA{} for the students, a difference of \esEndingLiftA{} \esEndingLiftACI{} standard deviations (\esEndingLiftB{} \esEndingLiftBCI{} when the model imitates a student). Overall positivity goes against the prediction for plain essays, which are less positive than the students' (\esNetA{} \esNetACI), and with it when the model imitates a student (\esNetB{} \esNetBCI). HL2 and HL3 are undecided for essays.

Two differences recur in all three comparisons, and I had predicted neither. Machine text puts less of its charge directly after an intensifier (\hcIntensified, \esIntensifiedA{} and \esIntensifiedB{} standard deviations) and less of it under a negation (\hcNegated, \esNegatedA{} and \esNegatedB, the last with an interval that reaches zero). People write \emph{very} and \emph{not} next to their feeling words more than these models do.

The differences are real and modest. A logistic regression on the ten measures tells human from machine text with an area under the curve of \hcAuc{} in HC3, \esAucA{} \esAucACI{} for plain essays and \esAucB{} \esAucBCI{} for essays written as a student (\figref{fig:machine}b). Word count alone does better in HC3 (\hcAucWords), where answer lengths differ, and is near chance for the length-matched essays (\esAucWordsA). The regression fitted on HC3 scores \esHcA{} and \esHcB{} on the essays, so part of the pattern carries from ChatGPT answers of 2023 to Claude essays of 2026. None of this makes a detector. Purpose-built detectors use far richer evidence and still misfire, notably against non-native writers\cite{liang2023,weberwulff2023}.

\dialParagraphs

\section{Discussion}

The aim was a measure of sentiment for which granularity holds. The signature provides one: charge adds, position sums shift, and the signature of a whole is the merge of the signatures of its parts, exactly in integer arithmetic and to rounding error in floating point. Totals, density, a level with a standard error, a predicted rating and an arc are all read from the same nineteen numbers, and the scale at which to read them can be chosen after the text has been scanned.

No component of this is new. Averaged word norms\cite{dodds2011}, separate positive and negative counts\cite{tausczik2010,golder2011}, mergeable moment summaries\cite{cormode2009,pebay2008} and the decompositions of pooled values\cite{olley1996,kitagawa1955} have long histories in their own fields, and the description of mergeable position weights is a classical result\cite{anselone1964,szekelyhidi1982}. The paper adds their combination for sentiment, an implementation tested for exactness, and evidence on how position and neutral words enter human judgements.

Of the two hypotheses the work began with, one failed and one survived in altered form. Feeling is not front-loaded. The end of a phrase, a sentence or a review counts for more than its start, and three tests on untouched data supported that revised claim. Neutral words do dilute, as the averaging account of impression formation predicts\cite{anderson1965,nisbett1981}, but mainly when they trail the feeling. A weighted average with a neutral first impression and weights that rise toward the end would describe all of these results. The shipped readout has the neutral prior and the buffer weight, and its position terms lean the same way.

The signature is not the most accurate reader of a single text. Transformers agree with raters better (\meanDistil{} to \meanRoberta{} against \meanSgw), and where the judgement of one text matters they are the better tool. The case for the signature rests on other properties: about \speedDistil{} times the throughput per processor core, a result that does not change with the way a text is cut, and readouts that can be audited word by word. One sensible division of labour is to scan everything with signatures and to send a larger model only the spans that the scan marks.

Several limits qualify these results. Word norms rate words, not events, so densities of narrative text should be read against the text's own baseline, and irony, sarcasm and meaning carried by syntax beyond the four operators are missed. The lexicon is English and its ratings date from the early 2010s. The constants were fitted on six corpora in which reviews and social media dominate, and charges refined on ratings did not travel across genres. No human ratings of chapters or books were available; above the sentence, validity rests on reviews and on fairy tales, where averaged VADER scores tracked overall valence better than the signature did. The standard error of a level is not a prediction interval for a human rating. Human judgements of long experiences are known to weigh peaks and endings rather than sums\cite{fredrickson1993}; the signature keeps what an averaging or an end-weighted readout needs, and a readout of peaks would need signatures kept per part. The comparison of human and machine writing covers \esN{} student essays on eight assignments, four models from one family reached through an agent harness, and ChatGPT answers from 2023; its judge was a model, not a human reader. Analyses added after the first results are marked exploratory, and the confirmatory tests of recency, though committed in advance, come from the same kinds of text as the tests that suggested it.

Three additions would test the measure further: human ratings of chapters, so that readouts above the sentence can be validated directly; word norms in other languages; and the decayed sum that Proposition~2 allows, for a running level in streams of speech or messages.
